\documentclass[11pt,a4paper]{article}
\usepackage[margin=25mm]{geometry}
\usepackage{fontspec}
\setmainfont{TeX Gyre Pagella}
\usepackage{amsmath,amssymb,mathtools}
\usepackage{graphicx}
\usepackage{xcolor}
\usepackage[unicode,colorlinks=true,linkcolor=blue,urlcolor=blue]{hyperref}
\usepackage{bookmark}
\usepackage{enumitem}
\setlist{itemsep=3pt,topsep=5pt}
\setlength{\emergencystretch}{3em}
\newcommand{\tightlist}{\setlength{\itemsep}{0pt}\setlength{\parskip}{0pt}}
\newcommand{\passthrough}[1]{#1}
\newcommand{\pandocbounded}[1]{#1}
\makeatletter
\def\maxwidth{\ifdim\Gin@nat@width>\linewidth\linewidth\else\Gin@nat@width\fi}
\def\maxheight{\ifdim\Gin@nat@height>0.75\textheight0.75\textheight\else\Gin@nat@height\fi}
\makeatother
\setkeys{Gin}{width=\maxwidth,height=\maxheight,keepaspectratio}
\hypersetup{pdfauthor={OpenAI GPT-6.1 Sol, Codex, Ultra effort}}
\begin{document}
\section{Idèles in extensions and their
cohomology}\label{iduxe8les-in-extensions-and-their-cohomology}

\emph{Written by OpenAI GPT-6.1 Sol in Codex, Ultra effort, October
2026. Self-checked by the writing AI; no independent review is claimed.
Public domain (CC0).}

A field extension replaces a place by the places above it. Its idèle
norm multiplies the corresponding local norms. This lets us compute
cyclic cohomology one place at a time. The resulting quotient of the
idèle group is usually infinite; the principal idèles impose relations
between those local obstructions.

We use the definitions and restricted product topology in \emph{Idèles
and the idèle class group}, lesson 3 of \emph{Adèles and L-functions},
the Hilbert 90 proof in
\href{../hilberts-theorem-90-and-kummer-theory.html}{Hilbert's theorem
90 and Kummer theory}, and finite local reciprocity in
\href{../local-reciprocity-and-norm-groups.html}{Local reciprocity and
norm groups}. The cyclic Tate groups and Herbrand quotient are defined
in
\href{../cohomology-of-cyclic-groups-and-the-herbrand-quotient.html}{Cohomology
of cyclic groups and the Herbrand quotient}.

Let \(K\) be a global field and \(L/K\) a finite separable extension. A
global field means a number field or a finite extension of a
finite-field rational function field. For a place \(v\) of \(K\), write
\(w\mid v\) for its extensions to \(L\). The arguments below apply to
either kind of global field. There are no archimedean places in the
function-field case.

\textbf{Prerequisite proof availability.} The named results below
identify specific programme lessons. The
\href{https://kokunoyumeto.github.io/open-math-courses/courses/NT-CFT/proof-dependencies.html\#lesson-13}{prerequisite
record} distinguishes published proofs, supplied owner texts awaiting
publication, and missing full proofs. A record or external reference is
not a supplied proof; arguments using an unavailable prerequisite retain
that dependency.

\subsection{1. Embeddings and norms}\label{embeddings-and-norms}

Write \[
J_K=\prod_v'K_v^\times,\qquad C_K=J_K/K^\times,
\] where the restricted product uses \(U_v=\mathcal O_v^\times\) at
finite places. If \(S\) is a finite set containing the archimedean
places, let \(S_L\) be its inverse image and put \[
J_L^S=\prod_{w\in S_L}L_w^\times\times
                \prod_{w\notin S_L}U_w.
\tag{1}
\] These are open subgroups, and their union is \(J_L\).

We recall the local decomposition in a form that also fixes the norm: \[
L\otimes_K K_v\simeq\prod_{w\mid v}L_w,\qquad
\sum_{w\mid v}[L_w:K_v]=[L:K].
\tag{2}
\] Indeed, write \(L=K[T]/(f)\) by a primitive element, factor its
separable polynomial over \(K_v\), and apply the polynomial Chinese
remainder theorem. Its irreducible factors give the finite extensions of
\(K_v\). The image of \(K[T]/(f)\) is dense in each such field, since
\(K\) is dense in \(K_v\) and powers of its generator span that field.
Thus those factors are exactly the completions at the extensions of
\(v\), proving (2).

In the Galois case, \(L_w\) contains the images of all conjugates of the
primitive element, since they already belong to \(L\). Thus it splits
\(f\), and \(L_w/K_v\) is Galois. Any local automorphism sends that
primitive element to a global conjugate, so it preserves the dense
subfield \(L\) and restricts to an element of \(G\) preserving \(w\).
Conversely every such global element extends continuously. This
identifies the local group with \(G_w\). Finally every other local
factor has a root of \(f\), which is some global conjugate in \(L_w\);
the resulting embedding is the embedding of \(L\) twisted by that global
automorphism. Its induced place is a conjugate of \(w\). Hence \(G\)
acts transitively on the places above \(v\).

Only finitely many places ramify. One way to see the finiteness needed
here is to choose a primitive element integral after excluding finitely
many denominators. Its nonzero discriminant is a unit outside a finite
further set. At such a finite place the reduction of its polynomial is
separable; lifting its relatively prime factors gives unramified local
factors. For a number field use integer rings; for a function field use
the integral closure of a finite-field polynomial ring and include the
places over infinity. This argument proves the same finite
exceptional-set assertion in both cases.

\subsubsection{Proposition 13.1. Extension and norm
maps}\label{proposition-13.1.-extension-and-norm-maps}

The diagonal embeddings at the places above \(v\) define a continuous
injection \(J_K\hookrightarrow J_L\). The continuous norm is \[
(N_{L/K}x)_v=\prod_{w\mid v}N_{L_w/K_v}(x_w).
\tag{3}
\] It sends principal idèles to their ordinary field norms, is
transitive in towers, and sends an embedded \(a\in J_K\) to
\(a^{[L:K]}\). It induces maps \[
C_K\hookrightarrow C_L,\qquad N_{L/K}:C_L\longrightarrow C_K.
\tag{4}
\]

\textbf{Proof.} A local unit stays a unit on extension, and the norm of
a local unit is a unit. Thus both formulas give idèles and are
continuous on every restricted product chart. The diagonal map is
injective because each component embedding is injective.

In (2), multiplication by \((x_w)\) is the product of the multiplication
operators on the factors, so its determinant is (3). For a principal
element \(x\in L^\times\), this operator is the scalar extension of
multiplication by \(x\) on the \(K\)-vector space \(L\). Its determinant
is therefore \(N_{L/K}(x)\), independently of \(v\). Transitivity
follows from field-norm transitivity at the completions and regrouping
the places. For \(a_v\in K_v^\times\), the determinant is
\(a_v^{[L:K]}\) by (2).

To prove the first injection in (4), we need
\(J_K\cap L^\times=K^\times\). Embed \(L\) in a finite Galois extension
\(M/K\). A diagonal idèle from \(J_K\) is fixed by
\(\operatorname{Gal}(M/K)\). If it is the principal idèle of
\(x\in L^\times\), each automorphism therefore fixes \(x\) in every
completion and hence fixes \(x\) as an element of \(M\). Thus
\(x\in K^\times\). The norm map descends because of the
principal-element calculation. ∎

For Galois \(L/K\), the norm (3) is also the group-module norm
\(\prod_{g\in G}g(x)\). This follows either from the determinant or by
grouping the embeddings according to the place they induce.

\subsection{2. Galois descent for
classes}\label{galois-descent-for-classes}

Assume \(L/K\) Galois, with \(G=\operatorname{Gal}(L/K)\). Its action on
idèles is \[
(g x)_w=g(x_{g^{-1}w}).
\tag{5}
\] Here the last \(g\) denotes the induced isomorphism of completions.

\subsubsection{Proposition 13.2. Fixed idèles and fixed
classes}\label{proposition-13.2.-fixed-iduxe8les-and-fixed-classes}

With the embeddings just defined, \[
J_L^G=J_K,\qquad C_L^G=C_K.
\tag{6}
\]

\textbf{Proof.} Fix \(w\mid v\) and its decomposition group \(G_w\). The
local field \(L_w/K_v\) is Galois with group \(G_w\); this follows also
from (2), since the global automorphisms permute its factors
transitively and the stabilizer acts on one factor. A fixed idèle has
\(x_w\in L_w^{G_w}=K_v\). Transitivity on the places above \(v\) then
makes all their components the image of the same element of
\(K_v^\times\). Almost all of these elements are units, proving the
first equality.

A fixed class need not initially have a fixed representative. If
\([x]\in C_L^G\), define \[
a_g=\frac{g(x)}x\in L^\times.
\] These principal elements satisfy \(a_{gh}=a_g\,g(a_h)\). The full,
noncyclic Hilbert 90 theorem proved in lesson 3 gives \(a_g=g(b)/b\) for
some \(b\in L^\times\). Consequently \(x/b\) is fixed by every \(g\),
hence belongs to \(J_K\), and represents \([x]\). Proposition 13.1 gives
injectivity, proving the second equality. ∎

This is the invariant part of the exact sequence
\(1\to L^\times\to J_L\to C_L\to1\). Hilbert 90 supplies the
surjectivity that invariants alone would not preserve.

\subsection{3. A cyclic induction
calculation}\label{a-cyclic-induction-calculation}

Suppose \(G=\langle\sigma\rangle\) has order \(n\). Let
\(H=\langle\tau\rangle\), \(\tau=\sigma^r\), be a subgroup, with
\(r=[G:H]\). We use additive notation for an abelian \(H\)-module \(M\).
Its induced module is \(B=M^r\), with \[
\sigma(x_0,\ldots,x_{r-1})=(\tau x_{r-1},x_0,\ldots,x_{r-2}).
\tag{7}
\]

\textbf{Lemma.} For \(i=0,-1\), \[
\widehat H^i(G,B)\simeq\widehat H^i(H,M).
\tag{8}
\]

\textbf{Proof.} Invariants in (7) are diagonal vectors \((a,\ldots,a)\)
with \(a\in M^H\). If \(N_H=1+\tau+\cdots+\tau^{n/r-1}\), then \[
N_G(x_0,\ldots,x_{r-1})
=(N_H\textstyle\sum x_i,\ldots,N_H\sum x_i).
\tag{9}
\] Thus invariants modulo norms give (8) in degree zero.

Summing coordinates modulo \((\tau-1)M\) identifies \[
B/(\sigma-1)B\simeq M/(\tau-1)M.
\tag{10}
\] To check the kernel as well as surjectivity, in the left quotient
each coordinate vector is equivalent to a vector in the first
coordinate, by successive applications of \(\sigma\). Returning after
\(r\) steps identifies its value with its image under \(\tau\). These
are precisely the relations in the right quotient; placing a value in
the first coordinate gives the inverse. Under (10), the norm to
invariants is exactly \(N_H\), by (9). Its kernel is
\(\widehat H^{-1}\), proving the other case. ∎

For a chosen \(w\mid v\), the identifications
\(L_w\simeq L_{\sigma^i w}\) identify \[
B_v=\prod_{w'\mid v}L_{w'}^\times
 \simeq\operatorname{Ind}_{G_w}^G L_w^\times.
\tag{11}
\] The same holds for the unit groups at finite places. This is (7)
written multiplicatively, not a trivial permutation action on the
factors.

\subsection{4. The restricted product contributes a direct
sum}\label{the-restricted-product-contributes-a-direct-sum}

For a cyclic unramified local extension, both Tate groups of its units
vanish. The unit norm is surjective by the unramified norm theorem used
in lesson 6. For the other group, Hilbert 90 writes a norm-one unit as
\(\tau(b)/b\). Since an unramified extension has the same valuation
group and a ground-field uniformizer, divide \(b\) by the ground-field
uniformizer to the power \(v(b)\). This makes \(b\) a unit without
changing its coboundary. Thus \[
\widehat H^0(G_w,U_w)=\widehat H^{-1}(G_w,U_w)=0
\quad\text{at an unramified finite place.}
\tag{12}
\]

\subsubsection{Proposition 13.3. Local decomposition of idèle
cohomology}\label{proposition-13.3.-local-decomposition-of-iduxe8le-cohomology}

For cyclic \(L/K\) and \(i=0,-1\), restriction to the local factors
gives \[
\widehat H^i(G,J_L)
\simeq\bigoplus_v\widehat H^i(G_w,L_w^\times).
\tag{13}
\] In particular, \[
J_K/N_{L/K}J_L
\simeq\bigoplus_v K_v^\times/N_{L_w/K_v}L_w^\times,
\qquad
\widehat H^{-1}(G,J_L)=0.
\tag{14}
\] One place \(w\) above each \(v\) is chosen; changing it transports
the same local quotient by a \(K_v\)-isomorphism.

\textbf{Proof.} Take finite \(S\) containing the archimedean and
ramified places. Invariants, the norm kernel and the image of
\(\sigma-1\) commute with arbitrary products of modules: the assertions
are componentwise, and for the image choose one preimage in each
component. Hence (8), (11) and (12) give \[
\widehat H^i(G,J_L^S)
=\prod_{v\in S}\widehat H^i(G_w,L_w^\times).
\tag{15}
\] There is no tail contribution. Passing to the union over finite \(S\)
changes this finite product into a direct sum. More explicitly, a
norm-zero or fixed representative belongs to some \(J_L^S\), and a
relation exhibiting a norm or coboundary also belongs to some such
chart; only finitely many group operations occur. Thus kernels and the
relevant images pass to the filtered union, proving (13).

Degree zero is invariants modulo norms, and (6) identifies invariants
with \(J_K\). Local degree zero is \(K_v^\times/N L_w^\times\), proving
the first formula in (14). The local degree-minus-one group is zero by
Hilbert 90, including at archimedean fields. This proves the last
assertion. ∎

There is a concrete interpretation of the first formula. A fixed idèle
is a unit at almost all unramified places, and those units are local
norms. Thus it has only finitely many nonzero obstruction classes.
Conversely any finite list of local classes can be represented by an
idèle. If all its local classes vanish, choose norm preimages that are
units outside a finite set, using (12); these preimages form an idèle
whose norm is the given one.

\subsection{5. The quotient for a finite set of
places}\label{the-quotient-for-a-finite-set-of-places}

\subsubsection{\texorpdfstring{Proposition 13.4. The Herbrand quotient
of
\(S\)-idèles}{Proposition 13.4. The Herbrand quotient of S-idèles}}\label{proposition-13.4.-the-herbrand-quotient-of-s-iduxe8les}

For cyclic \(L/K\), with \(S\) finite and containing all archimedean and
ramified places, \[
h(G,J_L^S)=\prod_{v\in S}[L_w:K_v].
\tag{16}
\]

\textbf{Proof.} Formula (15) makes both Tate groups finite. For a finite
cyclic extension of nonarchimedean local fields, finite reciprocity
gives norm index \([L_w:K_v]\), and Hilbert 90 gives degree-minus-one
group zero. Their Herbrand quotient is therefore \([L_w:K_v]\). At a
real place that becomes complex, the norm image is \(\mathbf R_{>0}\),
of index 2, and a norm-one complex number is \(z/\bar z\); hence the
same quotient is 2. Every other archimedean local degree is 1 and
contributes 1. Multiply these local orders in (15), with no tail
contribution by (12). ∎

The quotient \(h(G,J_L)\) need not be defined: its degree-zero Tate
group may be infinite. Formula (16) concerns the specified open
subgroup, not the whole idèle group.

\subsection{6. An infinite local quotient and a finite class
quotient}\label{an-infinite-local-quotient-and-a-finite-class-quotient}

Consider \(L=\mathbf Q(i)\), \(K=\mathbf Q\). At an odd prime \(p\), the
polynomial \(T^2+1\) splits if \(p\equiv1\pmod4\), and otherwise gives
the unramified quadratic extension. This follows from cyclicity of the
residue-field units and the simple-root lifting theorem. At 2 its
completion is the ramified quadratic field \(\mathbf Q_2(i)\). The real
completion becomes complex.

The norm groups proved in lessons 6 and 10 therefore give \[
J_{\mathbf Q}/N J_{\mathbf Q(i)}
\simeq \mathbf Z/2\ \oplus\ \mathbf Z/2\
        \oplus\!\!\bigoplus_{p\equiv3\ (4)}\mathbf Z/2.
\tag{17}
\] The first coordinate is the real sign. The second sends
\(x_2=2^\alpha u\) to the class of its odd unit \(u\) modulo
\(1+4\mathbf Z_2\). The \(p\)-coordinate is \(v_p(x_p)\) modulo 2. Split
primes contribute zero.

There are infinitely many primes \(3\pmod4\): from a putative finite
list take \(4\prod p-1\). It is \(3\pmod4\), none of the listed primes
divides it, and some prime factor is \(3\pmod4\), since a product of
primes all \(1\pmod4\) cannot be \(3\pmod4\). Thus (17) is infinite.

We can also compute the class quotient here without assuming global
reciprocity. Let \(e_\infty,e_2,e_p\) be the coordinate generators in
(17). The image of the principal idèle \(-1\) is \(e_\infty+e_2\); that
of a prime \(p\equiv3\pmod4\) is \(e_2+e_p\). Principal primes
\(1\pmod4\) and the prime 2 have zero image. These elements generate the
full principal image, since \(\mathbf Q^\times\) is generated by \(-1\)
and its primes. The quotient by these relations identifies every
coordinate generator with \(e_2\), and the sum of all coordinates
remains a nonzero invariant. Consequently \[
C_{\mathbf Q}/N C_{\mathbf Q(i)}\simeq\mathbf Z/2.
\tag{18}
\] Every sum here has finite support. The passage from (17) to (18)
displays the relations imposed by principal idèles.

\subsection{7. Exercises with solutions}\label{exercises-with-solutions}

\subsubsection{Exercise 1. The norm image is open ---
easy}\label{exercise-1.-the-norm-image-is-open-easy}

For finite separable \(L/K\), prove that \(N J_L\) is open in \(J_K\).

\textbf{Solution.} A norm subgroup for any finite extension of
nonarchimedean local fields is open: embed it in a finite Galois
extension, whose norm subgroup is open by local reciprocity, and use
norm transitivity to put that open subgroup inside the original norm
image. At infinity a norm image is the whole local multiplicative group
or the positive real subgroup; both are open. Choose finite \(S\)
containing ramification and infinity. At each \(v\in S\), the product of
the local norm images is an open subgroup \(V_v\subset K_v^\times\).
Outside \(S\), every unit is a norm of a unit from any one unramified
factor. Thus \[
\prod_{v\in S}V_v\times\prod_{v\notin S}U_v
\subset N J_L.
\] Choose the corresponding preimages componentwise; their unit tails
make them idèles. The displayed group is an open identity neighborhood,
so the subgroup \(N J_L\) is open. Infinite index is compatible with
openness.

\subsubsection{Exercise 2. Fixed classes ---
medium}\label{exercise-2.-fixed-classes-medium}

Prove that every \(G\)-fixed idèle class has a fixed representative.
Identify where Hilbert 90 enters.

\textbf{Solution.} Choose representative \(x\). Fixedness means
\(g(x)/x=a_g\in L^\times\), and the action gives \(a_{gh}=a_g g(a_h)\).
Hilbert 90 for the full finite Galois group gives \(a_g=g(b)/b\). Then
\(g(x/b)=x/b\), so the representative is in \(J_L^G=J_K\). Its class is
unchanged. Injection follows from \(J_K\cap L^\times=K^\times\). The
cocycle could not be removed merely by applying an exactness claim to
invariants; Hilbert 90 is precisely the missing surjectivity step.

\subsubsection{Exercise 3. The quadratic rational example ---
medium}\label{exercise-3.-the-quadratic-rational-example-medium}

Describe the quotient in (17), provide a representative for each
coordinate generator, and derive (18).

\textbf{Solution.} For \(e_\infty\), use real component \(-1\) and all
finite components 1. For \(e_2\), use component 3 at 2 and all other
components 1. For \(e_p\), use a uniformizer \(p\) at that inert prime
and 1 elsewhere. Their local norms fail exactly at the indicated place.
Local unramified norms have even valuation and arbitrary unit; the norm
group at 2 is \(2^{\mathbf Z}(1+4\mathbf Z_2)\); the real norm is
positive; at split primes it is surjective. This proves the coordinate
description, its independence and its exhaustiveness by (14). On
principal elements, \(-1\) gives \(e_\infty+e_2\), each inert prime
gives \(e_2+e_p\), and the remaining prime generators give zero.
Quotienting by their span identifies all generators, leaving a single
group of order 2 detected by the sum of coordinates. This is the class
quotient because \(N C_L\) is the image of \(N J_L\) after quotienting
by \(\mathbf Q^\times\).

\subsubsection{Exercise 4. Proving the decomposition directly ---
hard}\label{exercise-4.-proving-the-decomposition-directly-hard}

Prove both degrees of (13), keeping track of restricted-product support
and the twisted action on a block.

\textbf{Solution.} For a place block choose the cosets
\(1,\sigma,\ldots,\sigma^{r-1}\) of \(G_w=\langle\sigma^r\rangle\), and
identify the components using these automorphisms. The action is the
twisted cycle (7). Its invariants are diagonal \(G_w\)-invariants; its
norm is (9). Thus degree zero reduces to the local degree-zero quotient.
The sum map (10), with inverse given by the first coordinate, reduces
the kernel of the norm on coinvariants to
\(\ker N_{G_w}/(\sigma^r-1)L_w^\times\); this proves degree minus one,
too.

An idèle belongs to a chart with only finitely many unrestricted places.
At every remaining place the unit Tate groups vanish by unramified unit
norms and the unit-adjusted Hilbert 90 argument. For the product of
those tails, choose the local preimages all in unit blocks, so they form
a restricted-product preimage. Hence the chart's cohomology is the
finite product in (15). A representative or a finite algebraic relation
involves only one sufficiently large chart. Taking their union therefore
produces the direct sum over all places, not an unrestricted product of
local obstruction groups.

\subsection{Editable edition}\label{editable-edition}

The reading edition provides the complete LaTeX source of this lesson,
the cumulative course LaTeX and the editable source ZIP. The archive
contains all twenty-four Markdown lessons, complete LaTeX bodies,
original diagrams, metadata and reproduction instructions.

\subsection{References}\label{references}

Section 1 gives the local tensor decomposition from the primitive
element and polynomial remainder lemmas of lesson 1. Subsequent sections
prove norms, invariant class descent and the induced local calculation
in the restricted product.

\begin{itemize}
\tightlist
\item
  \href{https://www.mathi.uni-heidelberg.de/~schmidt/Neukirch-en/Neukirch_cft_02_may15.pdf}{Jürgen
  Neukirch, Class Field Theory --- The Bonn Lectures, Online Edition 2.0
  (May 2015), edited by Alexander Schmidt}.
\item
  \href{https://kskedlaya.org/cft/sec_abstractcft1.html}{Kiran S.
  Kedlaya, Notes on class field theory, author-hosted HTML edition}.
\end{itemize}

The \href{../free-proofs.html}{proof guide} gives the lesson sequence
and the exact prerequisite record. External references accompany the
written arguments.

\end{document}
